\chapter{The pipeline, and why it has that shape}
\label{ch:pipeline}

Binding a library with rosetta is a \emph{pipeline} rather than a single
command. That is not an implementation choice --- it is the only shape the
problem admits. This chapter explains why, because almost every surprise in
later chapters resolves back to it.

\section{The constraint everything follows from}

C++26 reflection answers questions \textbf{at the compile time of a program
that includes your headers}. It cannot be queried from a script, from a
configuration file, or from a build system. Only from C++, being compiled, with
your types in scope.

So to learn the shape of your types, something has to \textbf{compile}. To turn
those answers into files on disk, that something has to \textbf{run}.

\begin{center}
\emph{Compile, then run.} Everything below is that sequence, plus one step
before it and one after.
\end{center}

\section{The five steps}

\begin{center}
\begin{tikzpicture}[
  node distance=5mm,
  b/.style={draw,rounded corners=2pt,align=center,font=\scriptsize,inner sep=4pt,text width=23mm,minimum height=9mm},
  user/.style={b,fill=blue!6,draw=rosettablue},
  tool/.style={b,fill=orange!8,draw=rosettaorange},
  gen/.style={b,fill=black!4,draw=black!35,dashed},
  fin/.style={b,fill=green!7,draw=green!45!black},
  arr/.style={-{Stealth[length=2mm]},draw=black!55,font=\scriptsize}
]
\node[user] (m) {\texttt{manifest.json}\\+ your headers};
\node[tool,right=9mm of m] (rg) {\textbf{0}\ \texttt{rosetta\_gen}\\\scriptsize no reflection};
\node[gen,right=9mm of rg] (g) {\textbf{1}\ \texttt{gen/}\\driver source};
\node[tool,right=9mm of g] (drv) {\texttt{./generator}\\\scriptsize reflection HERE};
\node[gen,below=8mm of drv] (bind) {\textbf{2}\ \texttt{bindings/}\\one project per target};
\node[fin,left=9mm of bind] (built) {\textbf{3}\ compiled\\\texttt{.so .node .wasm}};
\node[fin,left=9mm of built] (load) {\textbf{4}\ \texttt{import}\\\texttt{require}};
\draw[arr] (m) -- (rg);
\draw[arr] (rg) -- (g);
\draw[arr] (g) -- node[above,align=center]{cmake\\build} (drv);
\draw[arr] (drv) -- (bind);
\draw[arr] (bind) -- node[above]{cmake} (built);
\draw[arr] (built) -- (load);
\end{tikzpicture}
\end{center}

\subsection*{Step 0 --- build \code{rosetta\_gen}}

\code{rosetta\_gen} is a JSON-to-C++ text templater. It does no reflection at
all and builds with an off-the-shelf C++17 compiler. You build it once.

It exists so that what \emph{you} write is a manifest and not a driver program.
Given \code{manifest.json} it emits three files: a \code{<name>.cpp} containing
a \code{main()} that reflects, a \code{bindings.h} with your types (and any
out-of-line annotations) baked in, and a \code{CMakeLists.txt} that knows where
the reflection toolchain lives.

\emph{Your library:} untouched. \code{rosetta\_gen} never opens your headers.
It only copies their names into the driver it writes.

\subsection*{Step 1 --- compile and run the driver}

Three commands, because ``compile then run'' is three things --- emit the
program, build it, execute it:

\begin{lstlisting}[style=sh]
rosetta_gen manifest.json gen     # emit the driver source
cmake -S gen -B gen/build         # COMPILE it -- reflection happens HERE
cmake --build gen/build
./generator bindings              # RUN it -- writes one project per target
\end{lstlisting}

During the \emph{compile}, clang-p2996 evaluates the reflection operators in
the driver: it walks your classes, enumerates fields and methods, reads
annotations, and resolves each member against the target language's type gate.
\textbf{Every decision about what can and cannot be bound is made at this
moment.} During the \emph{run}, the driver does nothing clever --- it prints
the conclusions it was compiled with.

Neither half can be skipped. A compiler cannot write your binding files; a
program cannot see types it was not compiled against.

\emph{Your library:} read for the \textbf{first} time, as source, by the C++26
toolchain.

\subsection*{Step 2 --- what lands in \code{bindings/}}

One self-contained project per target, plus \code{coverage.json} --- a
machine-readable account of what bound and what did not, available
\emph{before} anything is compiled.

The generated code is fully expanded (Section~\ref{sec:reflection-free}). No
splices, no \code{<experimental/meta>}, nothing that needs the fork.

\subsection*{Step 3 --- compile a backend}

Ordinary CMake, or npm, or emcmake, over an ordinary C++ project.

\emph{Your library:} read a \textbf{second} time --- the generated binding
\code{\#include}s your headers and calls your functions. This is where
\field{user\_sources} are compiled and \field{user\_lib} is linked in.

\subsection*{Step 4 --- load it}

\code{import}, \code{require}, \code{dlopen}. Nothing rosetta-specific is left.

\section{Your library is read three times, for different reasons}

\begin{center}
\begin{tabular}{@{}llll@{}}
\toprule
\textbf{Pass} & \textbf{By} & \textbf{To} & \textbf{Needs C++26} \\
\midrule
step 0 & \code{rosetta\_gen}       & \textbf{read} your comments  & \no \\
step 1 & the reflection driver  & \textbf{describe} your types & \yes \\
step 3 & the generated binding  & \textbf{call} them           & \no \\
\bottomrule
\end{tabular}
\end{center}

The middle row is the single most useful thing to hold in your head. The first
is the odd one out and worth a sentence: comments are not reflectable --- P2996
describes declarations, not the text around them --- so the documentation your
headers already carry is read \emph{lexically}, by the tool, while it is
emitting the driver. That pass compiles nothing and needs no toolchain; what it
finds is matched against what step~1 reflects, and dropped where the two
disagree. Chapter~\ref{ch:documentation} is about that pass.

The table has one further consequence worth stating on its own:

\begin{warn}
Annotations written \textbf{inline} --- \code{[[= rosetta::doc\{\ldots\} ]]} ---
pull \code{<experimental/meta>} into your header. Pass~2 then needs the C++26
toolchain as well, and the ``no reflection on the target'' promise is lost.
Out-of-line annotations (a \code{.ann.json} side-car) keep your header plain
C++ and pass~2 reflection-free. See Section~\ref{sec:outofline}.
\end{warn}

\section{Two commands, not five}

You rarely run the steps by hand. \code{rosetta\_gen --build} chains all of
them:

\begin{lstlisting}[style=sh]
rosetta_gen --build manifest.json
\end{lstlisting}

It emits \code{gen/}, builds and runs the driver, and then compiles every
backend the manifest declares --- each with its own build shape (plain
\code{cmake}, \code{npm install} + \code{npm run build}, \code{emcmake}, a
second-stage \code{dotnet build} or \code{mvn package}). A backend whose
toolchain is missing is \emph{skipped with a note}, not an error, and every run
ends with a per-backend summary.

\code{rosetta\_gen --clean manifest.json} takes you back to sources. Chapter~\ref{ch:rosettagen}
is the full reference for both.

The manual sequence remains useful for exactly the reasons you would expect:
inspecting \code{bindings.h}, hacking on the emitted driver, or wiring the
stages into a build system of your own.

\section{When to re-run what}

The two-pass structure decides this, and getting it wrong is the most common
way to be confused by rosetta.

\begin{center}
\small
\begin{tabular}{@{}L{6.2cm}L{7.4cm}@{}}
\toprule
\textbf{What changed} & \textbf{What to do} \\
\midrule
A field or method: added, removed, retyped &
  Rebuild the backend. The generated binding is expanded, so this needs a
  regeneration \emph{if the member set changed} --- re-run \code{--build}. \\
\addlinespace
An \textbf{inline} annotation &
  Re-run \code{--build}: the annotation is read during pass~1. \\
\addlinespace
A \textbf{side-car} \code{.ann.json} &
  Re-run \code{rosetta\_gen} \emph{and} the generator, not just the binding
  build --- the JSON is baked into \code{bindings.h} at \code{rosetta\_gen}
  time. \\
\addlinespace
A class added, or \field{targets} changed &
  Re-run the whole pipeline. \\
\addlinespace
A \field{user\_sources} glob now matches new files &
  Re-run \code{rosetta\_gen}: globs expand at generation time. \\
\addlinespace
The manifest moved &
  Re-run. Every relative path in it resolves from the manifest's own directory. \\
\bottomrule
\end{tabular}
\end{center}

\begin{note}
Re-emitting is idempotent: unchanged files are not rewritten, so the next
\code{cmake --build} recompiles nothing. Running the pipeline again when you
are not sure is cheap.
\end{note}

\section{The footgun the tool warns about}

If the emitted sources changed but a \code{generator} binary from a previous
emit still sits next to the output directory, running it without rebuilding
\textbf{silently re-emits the old bindings}. The tool detects this and prints a
rebuild note. Believe it.

The related one: a fresh checkout (or a \code{--clean}) leaves
\code{out\_dir/build} unconfigured, so \code{cmake --build} alone fails. The
\emph{configure} step is printed as a \code{next:} hint when that happens.
